\chapter{TPI for finite cardinals}

	\section{READ THIS}
	
		\noindent
		This comment is totally optional:\\
		\textless ** \textit{See unnecesary comment 0015, about stupid comments I've suffered. READ CAREFULLY PLEASE} \textgreater
		\\\\
		
		\noindent
		THIS ONE IS OBLIGATORY. Ignore the tittle (it is very short): \\
		\textless ** \textit{See unnecesary comment 0016, totally boring comment} \textgreater
		\\\\
		
		
	
	\section{The core idea through an example}
	
		\noindent
		Imagine two tables.
		\\\\
		
		\noindent
		In one table we have a little leather bag (LLB), labeled as "50". When we look inside, it has 50 gold coins inside. Over THIS table there is a box labeled as "DOMAIN", too.
		\\\\
		
		\noindent
		In the other table we have 50 LLBs. Each one labeled as 1,2,3,4,5,6,..., 48,49,50. We have not a singular LLB, one inside the other. We check them all, and each one has inside the same quantity of gold coins than the number that is labeled on it. And we have a box labeled as "ORIGIN".
		\\\\
		
		\noindent
		In each table, we empty all the bags, over THAT table, inside the box that is over each table.
		\\\\
		
		\noindent
		Question:\\\\
		There are more gold coins inside the box labeled as "DOMAIN" than inside the box labeled as "ORIGIN"?\\\\
		I hope your answer were clearly "NO".
		\\\\
		
		\noindent
		Inside the box "DOMAIN" we have 50 coins, and inside the box "ORIGIN" we have...\\
		$$ (1+50)*(50/2) =1275 $$\\
		MORE THAN THE DOUBLE !! Much more than the double... 1275 is WAY MORE than only the double of 50.
		\\\\
		
		\noindent
		We could extract gold coins, from the "ORIGIN" box, to refill each one of the 50 LLBs, with the proper quantity of gold coins, to have the same than the number that is labeled in each one.
		\\\\
		
		\noindent
		To create, a group of growing disjoint subsets, of gold coins, from an initial small state (1 coin inside the bag labeled as "1"), to a cardinal equal to the set of coins, in the other table ( 50 coins inside the bag labeled as "50", in THIS table).. we need more than the double of coins inside the original set, TO CREATE ALL THE SUBSETS: a set with 1275 coins.
		\\\\
		
		\noindent
		We are going to create THIS phenomenom (TPI), adapted to infinity and two different alephs. Having $\mathcal{P(\mathbb{N})}$ as the DOMAIN and $\mathbb{N}$ as the ORIGIN set.
		\\\\
		
		\noindent
		We call it "ORIGIN" because it is the origin from where we extract each disjoint subset.
		\\\\
		
		\noindent
		Do you think it is "impossible"? See? My second "mathematically impossible" intermediate point.
		\\\\
	
	\section {The same example, but for groups instead of element by element}
	
		\noindent
		REREAD the previous example, but adding a letter "K" to each quantity inside each LLB.
		\\\\
		
		\noindent
		In the DOMAIN table, we have a LLB with 50K golds coins inside: 50 000.
		\\\\
		
		\noindent
		Over the table ORIGIN, that has the box ORIGIN, we have 50 different LLBs, labeled as 1K, 2K, 3K, ..., 48K, 49K and 50K, each one with the proper quantity of gold coins, inside, respecting their labels: 1000, 2000, 3000, ..., 48000, 49000, 50000. Inside the box ORIGIN, we will have 1 275 000 gold coins: the same proportion than in the previous example, but one million gold coins more (more or less) instead only one thousand.
		\\\\
		
		\noindent
		I one example we build the subsets, incrementing the cardinal one element each time. 50 steps in total. In the second example, we had 50 steps, but much more coins involved in each step. We grow in one by singular elements, and in another one by groups. The observation is the same, growing by elements than growing by groups.
		\\\\
		
		\noindent
		Now imagine that in the table DOMAIN, in its LLB we have 50000 gold coins... but we want to create the subsets growing element by element. We will need intead of 50 steps, 50000 steps... and:\\
		$$ (1 + 50000)* (50000/2) = 1250025000 \:\: gold \:\: coins $$
		\\\\
		
		\noindent
		The more steps we had, the greater needs to be the set ORIGIN, to being able to create ALL disjoint subsets of each step.
		\\\\
		
		\noindent
		The more steps we had, the more obvious is the initial observation. The proportion of the difference, bewteen ORIGIN and the original DOMAIN set, only grows more and more in a ridicoulous way. IN FINITE SPECTRUM.
		\\\\
		
		\noindent
		Creating a TPI between an ORIGIN set and a DOMAIN set, easily, means a case MUCH WORST than doubling the elements inside DOMAIN.
		\\\\
		
		\noindent
		Imagine a dude saying: we have the set X with cardinal $\aleph_{k}$. Now we double its elements, creating a prima-copy of each element inside X. We will call the set with normal-elements and theirs prima-copies: $X'$. Imagine this dude add: like the Cantor's Theorem is proven, WE KNOW $X'$ has a cardinal much smaller than $\aleph_{k}$. EVEN after doubling their elements correctly. What would you say to that dude? EXACTLY... $<$ Enter very bad words here $>$.
		\\\\
		
		\noindent
		And I repeat, just a TPI, with a few steps, with a few subsets created... means easily a case WORST than doubling the elements inside the DOMAIN, inside the ORIGIN set. We need to adapt the concept "growing subsets" between two consecutive alephs without any other cardinal betweeen ( forget Continuum Hypothesis ). In finite spectrum we have many different finite cardinals; but between $\aleph_{0}$ and $\aleph_{1}$... we haven't got that easy tool.
		\\\\
		
		\noindent
		We need to be sure, the steps grows until the DOMAIN set size, but we can not use the word "cardinal". HOW?? We need two concepts: the semi-xmas-tree pattern to measure proximity; and the group of concepts, NWSP, WSP, QRE, to prove "improvements" between two consecutive alephs. This last one will let us to define the "steps" clearly.
		\\\\
	
	
	\newpage	
	\section{SEMI-XMAS-TREE pattern in finite cardinalities:}
	
		\noindent
		
		
		
		
		
		
		
		
		
		